%! Author = nursace
%! Date = 16.09.2026
\documentclass[aspectratio=169,11pt]{beamer}
% Week 1 Practice — Setup, first notebook, predict-by-hand
% Course: AI Frameworks (ages 16--18)

% Shared Beamer preamble for AI Frameworks
% Audience: college students ages 16--18 (intuition over math)
% Include with: % Shared Beamer preamble for AI Frameworks
% Audience: college students ages 16--18 (intuition over math)
% Include with: % Shared Beamer preamble for AI Frameworks
% Audience: college students ages 16--18 (intuition over math)
% Include with: \input{preamble}

\usetheme{Madrid}
\usecolortheme{default}

% Clean palette: deep slate + teal accent (distinct from Java navy decks)
\definecolor{LectureNavy}{HTML}{1A365D}
\definecolor{LectureAccent}{HTML}{319795}
\definecolor{LectureGray}{HTML}{4A5568}
\definecolor{CodeBg}{HTML}{F7FAFC}
\definecolor{AlertBg}{HTML}{FFF5F5}
\definecolor{TryBg}{HTML}{F0FFF4}
\definecolor{QuizBg}{HTML}{EBF8FF}
\definecolor{AnswerKeyBg}{HTML}{FFFFF0}
\definecolor{AnswerGold}{HTML}{B7791F}
\definecolor{KeywordBlue}{HTML}{2B6CB0}
\definecolor{StringGreen}{HTML}{276749}
\definecolor{CommentGray}{HTML}{718096}

\setbeamercolor{structure}{fg=LectureNavy}
\setbeamercolor{palette primary}{bg=LectureNavy,fg=white}
\setbeamercolor{palette secondary}{bg=LectureAccent,fg=white}
\setbeamercolor{palette tertiary}{bg=LectureGray,fg=white}
% Madrid title box is dark --- title/subtitle must be light for contrast
\setbeamercolor{title}{fg=white,bg=LectureNavy}
\setbeamercolor{subtitle}{fg=white,bg=LectureNavy}
\setbeamercolor{author}{fg=LectureNavy}
\setbeamercolor{date}{fg=LectureGray}
\setbeamercolor{institute}{fg=LectureGray}
\setbeamercolor{frametitle}{bg=LectureNavy,fg=white}
\setbeamercolor{block title}{bg=LectureNavy,fg=white}
\setbeamercolor{block body}{bg=CodeBg,fg=black}
\setbeamercolor{block title alerted}{bg=LectureAccent,fg=white}
\setbeamercolor{block body alerted}{bg=TryBg,fg=black}
\setbeamercolor{block title example}{bg=LectureGray,fg=white}
\setbeamercolor{block body example}{bg=AlertBg,fg=black}

\setbeamertemplate{navigation symbols}{}
\setbeamertemplate{itemize items}[circle]
\setbeamertemplate{enumerate items}[default]

\usepackage[T1]{fontenc}
\usepackage[utf8]{inputenc}
\usepackage{lmodern}
\usepackage{booktabs}
\usepackage{tabularx}
\usepackage{graphicx}
\usepackage{hyperref}
\usepackage{listings}
\usepackage{xcolor}
\usepackage{tikz}
\usetikzlibrary{positioning,arrows.meta,shapes.geometric}

% listings: Python without shell-escape
\lstdefinestyle{pythonstyle}{
    language=Python,
    basicstyle=\ttfamily\small,
    keywordstyle=\color{KeywordBlue}\bfseries,
    commentstyle=\color{CommentGray}\itshape,
    stringstyle=\color{StringGreen},
    numbers=left,
    numberstyle=\tiny\color{CommentGray},
    numbersep=6pt,
    frame=single,
    rulecolor=\color{LectureGray!40},
    backgroundcolor=\color{CodeBg},
    breaklines=true,
    showstringspaces=false,
    tabsize=4,
    captionpos=b,
}
\lstset{style=pythonstyle}

\newcommand{\pycode}[1]{{\ttfamily\small\detokenize{#1}}}
\newcommand{\trythis}[1]{%
    \begin{alertblock}{Try this}
    #1
    \end{alertblock}%
}
\newcommand{\commonmistake}[1]{%
    \begin{exampleblock}{Common mistake}
    #1
    \end{exampleblock}%
}
\newcommand{\remember}[1]{%
    \begin{block}{Remember}
    #1
    \end{block}%
}

% Formative in-class checks: question slide, then a dedicated Answer slide
\newcommand{\quizpause}{%
    \par\vspace{0.5em}
    {\small\textit{Pause --- answer (clicker / raise hand / neighbour) before the next slide.}}%
}
\newcommand{\quizbox}[1]{%
    \setbeamercolor{block title}{bg=KeywordBlue,fg=white}%
    \setbeamercolor{block body}{bg=QuizBg,fg=black}%
    \begin{block}{Quiz}
    #1
    \end{block}%
    \setbeamercolor{block title}{bg=LectureNavy,fg=white}%
    \setbeamercolor{block body}{bg=CodeBg,fg=black}%
}
\newcommand{\answerbox}[1]{%
    \setbeamercolor{block title}{bg=AnswerGold,fg=white}%
    \setbeamercolor{block body}{bg=AnswerKeyBg,fg=black}%
    \begin{block}{Answer}
    #1
    \end{block}%
    \setbeamercolor{block title}{bg=LectureNavy,fg=white}%
    \setbeamercolor{block body}{bg=CodeBg,fg=black}%
}

% Empty author/date placeholders for the lecturer to fill
\author{}
\date{}


\usetheme{Madrid}
\usecolortheme{default}

% Clean palette: deep slate + teal accent (distinct from Java navy decks)
\definecolor{LectureNavy}{HTML}{1A365D}
\definecolor{LectureAccent}{HTML}{319795}
\definecolor{LectureGray}{HTML}{4A5568}
\definecolor{CodeBg}{HTML}{F7FAFC}
\definecolor{AlertBg}{HTML}{FFF5F5}
\definecolor{TryBg}{HTML}{F0FFF4}
\definecolor{QuizBg}{HTML}{EBF8FF}
\definecolor{AnswerKeyBg}{HTML}{FFFFF0}
\definecolor{AnswerGold}{HTML}{B7791F}
\definecolor{KeywordBlue}{HTML}{2B6CB0}
\definecolor{StringGreen}{HTML}{276749}
\definecolor{CommentGray}{HTML}{718096}

\setbeamercolor{structure}{fg=LectureNavy}
\setbeamercolor{palette primary}{bg=LectureNavy,fg=white}
\setbeamercolor{palette secondary}{bg=LectureAccent,fg=white}
\setbeamercolor{palette tertiary}{bg=LectureGray,fg=white}
% Madrid title box is dark --- title/subtitle must be light for contrast
\setbeamercolor{title}{fg=white,bg=LectureNavy}
\setbeamercolor{subtitle}{fg=white,bg=LectureNavy}
\setbeamercolor{author}{fg=LectureNavy}
\setbeamercolor{date}{fg=LectureGray}
\setbeamercolor{institute}{fg=LectureGray}
\setbeamercolor{frametitle}{bg=LectureNavy,fg=white}
\setbeamercolor{block title}{bg=LectureNavy,fg=white}
\setbeamercolor{block body}{bg=CodeBg,fg=black}
\setbeamercolor{block title alerted}{bg=LectureAccent,fg=white}
\setbeamercolor{block body alerted}{bg=TryBg,fg=black}
\setbeamercolor{block title example}{bg=LectureGray,fg=white}
\setbeamercolor{block body example}{bg=AlertBg,fg=black}

\setbeamertemplate{navigation symbols}{}
\setbeamertemplate{itemize items}[circle]
\setbeamertemplate{enumerate items}[default]

\usepackage[T1]{fontenc}
\usepackage[utf8]{inputenc}
\usepackage{lmodern}
\usepackage{booktabs}
\usepackage{tabularx}
\usepackage{graphicx}
\usepackage{hyperref}
\usepackage{listings}
\usepackage{xcolor}
\usepackage{tikz}
\usetikzlibrary{positioning,arrows.meta,shapes.geometric}

% listings: Python without shell-escape
\lstdefinestyle{pythonstyle}{
    language=Python,
    basicstyle=\ttfamily\small,
    keywordstyle=\color{KeywordBlue}\bfseries,
    commentstyle=\color{CommentGray}\itshape,
    stringstyle=\color{StringGreen},
    numbers=left,
    numberstyle=\tiny\color{CommentGray},
    numbersep=6pt,
    frame=single,
    rulecolor=\color{LectureGray!40},
    backgroundcolor=\color{CodeBg},
    breaklines=true,
    showstringspaces=false,
    tabsize=4,
    captionpos=b,
}
\lstset{style=pythonstyle}

\newcommand{\pycode}[1]{{\ttfamily\small\detokenize{#1}}}
\newcommand{\trythis}[1]{%
    \begin{alertblock}{Try this}
    #1
    \end{alertblock}%
}
\newcommand{\commonmistake}[1]{%
    \begin{exampleblock}{Common mistake}
    #1
    \end{exampleblock}%
}
\newcommand{\remember}[1]{%
    \begin{block}{Remember}
    #1
    \end{block}%
}

% Formative in-class checks: question slide, then a dedicated Answer slide
\newcommand{\quizpause}{%
    \par\vspace{0.5em}
    {\small\textit{Pause --- answer (clicker / raise hand / neighbour) before the next slide.}}%
}
\newcommand{\quizbox}[1]{%
    \setbeamercolor{block title}{bg=KeywordBlue,fg=white}%
    \setbeamercolor{block body}{bg=QuizBg,fg=black}%
    \begin{block}{Quiz}
    #1
    \end{block}%
    \setbeamercolor{block title}{bg=LectureNavy,fg=white}%
    \setbeamercolor{block body}{bg=CodeBg,fg=black}%
}
\newcommand{\answerbox}[1]{%
    \setbeamercolor{block title}{bg=AnswerGold,fg=white}%
    \setbeamercolor{block body}{bg=AnswerKeyBg,fg=black}%
    \begin{block}{Answer}
    #1
    \end{block}%
    \setbeamercolor{block title}{bg=LectureNavy,fg=white}%
    \setbeamercolor{block body}{bg=CodeBg,fg=black}%
}

% Empty author/date placeholders for the lecturer to fill
\author{}
\date{}


\usetheme{Madrid}
\usecolortheme{default}

% Clean palette: deep slate + teal accent (distinct from Java navy decks)
\definecolor{LectureNavy}{HTML}{1A365D}
\definecolor{LectureAccent}{HTML}{319795}
\definecolor{LectureGray}{HTML}{4A5568}
\definecolor{CodeBg}{HTML}{F7FAFC}
\definecolor{AlertBg}{HTML}{FFF5F5}
\definecolor{TryBg}{HTML}{F0FFF4}
\definecolor{QuizBg}{HTML}{EBF8FF}
\definecolor{AnswerKeyBg}{HTML}{FFFFF0}
\definecolor{AnswerGold}{HTML}{B7791F}
\definecolor{KeywordBlue}{HTML}{2B6CB0}
\definecolor{StringGreen}{HTML}{276749}
\definecolor{CommentGray}{HTML}{718096}

\setbeamercolor{structure}{fg=LectureNavy}
\setbeamercolor{palette primary}{bg=LectureNavy,fg=white}
\setbeamercolor{palette secondary}{bg=LectureAccent,fg=white}
\setbeamercolor{palette tertiary}{bg=LectureGray,fg=white}
% Madrid title box is dark --- title/subtitle must be light for contrast
\setbeamercolor{title}{fg=white,bg=LectureNavy}
\setbeamercolor{subtitle}{fg=white,bg=LectureNavy}
\setbeamercolor{author}{fg=LectureNavy}
\setbeamercolor{date}{fg=LectureGray}
\setbeamercolor{institute}{fg=LectureGray}
\setbeamercolor{frametitle}{bg=LectureNavy,fg=white}
\setbeamercolor{block title}{bg=LectureNavy,fg=white}
\setbeamercolor{block body}{bg=CodeBg,fg=black}
\setbeamercolor{block title alerted}{bg=LectureAccent,fg=white}
\setbeamercolor{block body alerted}{bg=TryBg,fg=black}
\setbeamercolor{block title example}{bg=LectureGray,fg=white}
\setbeamercolor{block body example}{bg=AlertBg,fg=black}

\setbeamertemplate{navigation symbols}{}
\setbeamertemplate{itemize items}[circle]
\setbeamertemplate{enumerate items}[default]

\usepackage[T1]{fontenc}
\usepackage[utf8]{inputenc}
\usepackage{lmodern}
\usepackage{booktabs}
\usepackage{tabularx}
\usepackage{graphicx}
\usepackage{hyperref}
\usepackage{listings}
\usepackage{xcolor}
\usepackage{tikz}
\usetikzlibrary{positioning,arrows.meta,shapes.geometric}

% listings: Python without shell-escape
\lstdefinestyle{pythonstyle}{
    language=Python,
    basicstyle=\ttfamily\small,
    keywordstyle=\color{KeywordBlue}\bfseries,
    commentstyle=\color{CommentGray}\itshape,
    stringstyle=\color{StringGreen},
    numbers=left,
    numberstyle=\tiny\color{CommentGray},
    numbersep=6pt,
    frame=single,
    rulecolor=\color{LectureGray!40},
    backgroundcolor=\color{CodeBg},
    breaklines=true,
    showstringspaces=false,
    tabsize=4,
    captionpos=b,
}
\lstset{style=pythonstyle}

\newcommand{\pycode}[1]{{\ttfamily\small\detokenize{#1}}}
\newcommand{\trythis}[1]{%
    \begin{alertblock}{Try this}
    #1
    \end{alertblock}%
}
\newcommand{\commonmistake}[1]{%
    \begin{exampleblock}{Common mistake}
    #1
    \end{exampleblock}%
}
\newcommand{\remember}[1]{%
    \begin{block}{Remember}
    #1
    \end{block}%
}

% Formative in-class checks: question slide, then a dedicated Answer slide
\newcommand{\quizpause}{%
    \par\vspace{0.5em}
    {\small\textit{Pause --- answer (clicker / raise hand / neighbour) before the next slide.}}%
}
\newcommand{\quizbox}[1]{%
    \setbeamercolor{block title}{bg=KeywordBlue,fg=white}%
    \setbeamercolor{block body}{bg=QuizBg,fg=black}%
    \begin{block}{Quiz}
    #1
    \end{block}%
    \setbeamercolor{block title}{bg=LectureNavy,fg=white}%
    \setbeamercolor{block body}{bg=CodeBg,fg=black}%
}
\newcommand{\answerbox}[1]{%
    \setbeamercolor{block title}{bg=AnswerGold,fg=white}%
    \setbeamercolor{block body}{bg=AnswerKeyBg,fg=black}%
    \begin{block}{Answer}
    #1
    \end{block}%
    \setbeamercolor{block title}{bg=LectureNavy,fg=white}%
    \setbeamercolor{block body}{bg=CodeBg,fg=black}%
}

% Empty author/date placeholders for the lecturer to fill
\author{}
\date{}


\title{Week 1 Practice: Setup \& First Predictions}
\subtitle{AI Frameworks}

\begin{document}

    \begin{frame}
        \titlepage
    \end{frame}

    \begin{frame}{Practice goals (80 minutes)}
        By the end of this session you should have:
        \begin{enumerate}
            \item A working environment (Colab \emph{or} local Jupyter)
            \item A notebook where you can run cells and see output
            \item A simple plot from a tiny list of numbers
            \item A hand-made rule that ``predicts'' a house price from size
        \end{enumerate}
    \end{frame}

    \begin{frame}{Agenda}
        \begin{enumerate}
            \item Environment check (15--20 min)
            \item Notebook warm-up (15 min)
            \item First tiny plot (15 min)
            \item Predict-by-hand activity (20--25 min)
            \item Done checklist \& share-out (10 min)
        \end{enumerate}
    \end{frame}

%======================================================================
    \section{Environment}
%======================================================================

    \begin{frame}{Option A --- Google Colab (recommended if unsure)}
        \begin{enumerate}
            \item Open \url{https://colab.research.google.com}
            \item Sign in with a school or personal Google account
            \item File $\rightarrow$ New notebook
            \item Run a cell: \pycode{print(1+1)} --- you should see \pycode{2}
        \end{enumerate}
        \vspace{0.5em}
        No local install needed; runs in the browser.
    \end{frame}

    \begin{frame}{Option B --- Local Jupyter}
        \begin{itemize}
            \item If the instructor prepared Anaconda / VS~Code: follow the handout
            \item Create a folder for this course
            \item Launch Jupyter Lab or Notebook; create \pycode{week01.ipynb}
        \end{itemize}
        \vspace{0.5em}
        \commonmistake{Spending the whole session only on installers.
        If stuck after 15 minutes, switch to Colab and continue learning.}
    \end{frame}

    \begin{frame}[fragile]{Checkpoint A --- verify Python}
        \begin{lstlisting}
import sys
print(sys.version)

import matplotlib
print("matplotlib OK")
        \end{lstlisting}
        \vspace{0.3em}
        If both print without errors, you are ready.
    \end{frame}

%======================================================================
    \section{Notebook warm-up}
%======================================================================

    \begin{frame}{How a notebook works}
        \begin{itemize}
            \item \textbf{Code cells:} run Python
            \item \textbf{Markdown cells:} notes and headings (optional today)
            \item Run order matters: if a cell uses a variable, define it first
        \end{itemize}
        \vspace{0.5em}
        \trythis{Restart the kernel once today on purpose, then re-run cells from the top.}
    \end{frame}

    \begin{frame}[fragile]{Warm-up cells}
        \begin{lstlisting}
print("Hello, AI Frameworks!")

name = "Alex"   # change me
print("Student:", name)

sizes = [50, 65, 80, 95]
print("Number of houses:", len(sizes))
print("First size:", sizes[0])
        \end{lstlisting}
    \end{frame}

%======================================================================
    \section{First plot}
%======================================================================

    \begin{frame}[fragile]{A minimal plot}
        \begin{lstlisting}
import matplotlib.pyplot as plt

sizes = [50, 65, 80, 95]
prices = [180, 200, 210, 250]  # made-up

plt.scatter(sizes, prices)
plt.xlabel("Size (m^2)")
plt.ylabel("Price (k)")
plt.title("Tiny housing example")
plt.show()
        \end{lstlisting}
    \end{frame}

    \begin{frame}{What to notice}
        \begin{itemize}
            \item Larger size tends to go with higher price \emph{in this toy data}
            \item Real data will be messier---that is normal
            \item Plots help you think before you model
        \end{itemize}
    \end{frame}

%======================================================================
    \section{Predict by hand}
%======================================================================

    \begin{frame}{Activity: invent a rule}
        Goal: predict \textbf{price} from \textbf{size} only---no ML library.
        \vspace{0.5em}
        \begin{enumerate}
            \item Look at the scatter plot (or the table of pairs)
            \item Invent a simple rule, e.g.\ \pycode{price = a * size + b}
            \item Pick numbers for \pycode{a} and \pycode{b} that feel reasonable
            \item Predict for a \emph{new} size (e.g.\ 70~m$^2$) not in the list
        \end{enumerate}
    \end{frame}

    \begin{frame}[fragile]{Template}
        \begin{lstlisting}
def predict_price(size):
    a = 2.0    # tweak these
    b = 50
    return a * size + b

print(predict_price(70))

# Optional: compare your rule to the toy points
for s, p in zip(sizes, prices):
    print(s, "real:", p, "my guess:", predict_price(s))
        \end{lstlisting}
    \end{frame}

    \begin{frame}{Discuss with a neighbour (3 minutes)}
        \begin{itemize}
            \item Whose rule fits the points better ``by eye''?
            \item What happens if you change \pycode{a} a lot?
            \item Why might two people choose different rules?
        \end{itemize}
        \vspace{0.5em}
        Next weeks: algorithms will search for good rules using data---same idea, more systematic.
    \end{frame}

%--- Short formative check (optional projectable quiz) ---
    \begin{frame}{Quiz --- Practice check}
        \quizbox{In supervised learning terms, house \textbf{size} is a \ldots{} and \textbf{price} is a \ldots}
        \begin{enumerate}
            \item label ; feature
            \item feature ; label
            \item library ; notebook
            \item plot ; colour
        \end{enumerate}
        \quizpause
    \end{frame}

    \begin{frame}{Answer --- Practice check}
        \answerbox{\textbf{2.} Size is a \textbf{feature} (input); price is the \textbf{label} / target.}
    \end{frame}

%======================================================================
    \section{Done checklist}
%======================================================================

    \begin{frame}{Done checklist}
        Tick when finished:
        \begin{enumerate}
            \item Environment works (Colab or local)
            \item Notebook runs \pycode{print} and imports \pycode{matplotlib}
            \item Scatter plot of size vs price appears
            \item \pycode{predict\_price} function returns a number for size 70
            \item You can say in one sentence what a feature and a label are
        \end{enumerate}
    \end{frame}

    \begin{frame}{Before you leave}
        \begin{itemize}
            \item Save / download your notebook (\pycode{week01.ipynb})
            \item Note one thing that was confusing---bring it to Week 2
        \end{itemize}
        \vspace{0.6em}
        \begin{block}{Next week}
            Lecture: Python enough for data (lists, dicts, loops, functions)\\
            Practice: small coding warm-ups without pandas (yet)
        \end{block}
    \end{frame}

\end{document}